\section{Background and Related Work}
\label{sec:related}

\paragraph{Locate-then-edit editing.} ROME \citep{meng2022rome} writes a
single rank-one update to one MLP layer; MEMIT \citep{meng2023memit} spreads a
batch of edits across several layers; AlphaEdit \citep{fang2025alphaedit}
applies the ROME/MEMIT target but first projects the update onto the
null-space of preserved-knowledge keys. ENCORE-style analyses
\citep{gupta2024encore} attribute damage to the growth of the edited matrix's
norm. \citet{hase2023doeslocalization} showed that where a fact is
\emph{localized} need not be where it is best \emph{edited}, motivating
mechanism-level rather than localization-level accounts. Lifelong-editing
failures under superposition \citep{superposition2024} further motivate a
per-edit interference predictor.

\paragraph{Hypernetwork and meta-learning editors.} A second family replaces
the closed-form solve with a learned update rule. KE \citep{decao2021ke} trains
a hypernetwork to predict a constrained parameter edit from the (subject,
relation, new object) triple; MEND \citep{mitchell2022mend} learns a low-rank
transform of the fine-tuning gradient so a single reshaped gradient step
installs the edit with bounded footprint; SERAC \citep{mitchell2022serac}
sidesteps editing the base weights, routing edited queries to a small auxiliary
scope classifier and counterfactual model. Because their update is a trained
network's output rather than a minimum-norm linear solve, these editors do not
admit the closed-form rank-one algebra of Section~\ref{sec:method}; a
key-cosine account of their damage would need a different derivation, and we
treat this as an explicit extension target (Section~\ref{sec:dissociation})
rather than assume the present law transfers.

\paragraph{Memory and in-context editors.} A third family avoids modifying
the base weights at inference time by storing edits in an external, queryable
structure. GRACE \citep{hartvigsen2023grace} attaches a discrete key-value
adaptor at a chosen layer and retrieves a stored value when an input's
activation falls within a learned neighborhood of an edited key; WISE
\citep{wang2024wise} extends this with a side-memory architecture designed to
stay stable under the long edit sequences our stress test probes in
Section~\ref{sec:sequential}. Making no in-place weight perturbation, these
methods admit no weight-level collateral interference; but their retrieval
geometry -- whether a probe's activation falls inside an edited key's
neighborhood -- is a natural analogue of the key-cosine coupling we study, and
the most direct route to a future geometric account of their failure modes.

\paragraph{Editing-quality evaluation.} The canonical score for a
locate-then-edit update decomposes into edit success (does the model now
answer the target fact), generalization to paraphrases, and locality on a
neighborhood of unrelated-but-related facts
\citep{meng2022rome,meng2023memit}; we adopt the same
Efficacy/Paraphrase/Neighborhood triple as the edit-quality complement to the
damage spectrum in Section~\ref{sec:dissociation}, so a locality-preserving
editor and a geometry-predicted-low-damage edit can be told apart from one
that simply fails to install the fact.

\paragraph{Gradient-influence and instance-attribution theory.} The
theoretical anchor for the \SxC{} surrogate is the gradient-influence line of
work, which attributes a model's prediction on one input to training updates
caused by another: influence functions \citep{koh2017influence} via an
implicit-Hessian approximation to leave-one-out retraining, and TracIn/GradSim
\citep{pruthi2020gradsim} via a first-order, checkpoint-averaged dot product of
gradients cheap enough to compute over a full edit set. Section~\ref{sec:method}
shows that for a rank-one update this first-order influence has an exact closed
form in the edit and probe keys alone, with no gradient computation -- the
\SxC{} decomposition is the GradSim family specialized to the rank-one update,
not an independent heuristic.

\paragraph{Positioning against CLaRE-ty.} CLaRE-ty
\citep{clarety2026} is the closest concurrent work: an a-priori per-fact
interference predictor from a single-layer forward-activation cosine
$\cos(h_i^L,h_j^L)$, evaluated over \clareNfacts{} facts, five locate-then-edit
editors, and three attention models, baselined only against GradSim, and
correlational by the authors' own Limitations. We are complementary rather
than competing -- the mechanistic, causal account behind that empirical
signal -- and differ on four axes, each still correct under the Llama-family
scoping:
(i)~\emph{mechanism vs.\ metric} -- a closed-form ROME \SxC{} decomposition
of \dW{} (Section~\ref{sec:method}), a zero-cost GradSim surrogate, where
CLaRE-ty measures an empirical hidden-state cosine and derives no algebra;
(ii)~\emph{norm-growth head-to-head and regime} -- key-cosine vs.\ ENCORE
norm-growth on a confound-clean metric, the L12$\rightarrow$L14 crossover,
and a sign-tracks-regime scale law, where CLaRE-ty's only baseline is
GradSim; (iii)~\emph{editor\slash architecture dissociation} -- ROME vs.\
fine-tuning (FT), Kullback--Leibler-regularized FT (KL-FT), and MEMIT, and
Llama vs.\ gemma/Phi/Qwen, all 3-seed, where CLaRE-ty evaluates neither
fine-tuning editors nor Qwen/Mistral predictively; and
(iv)~\emph{causal test} -- AlphaEdit null-space projection removes the
geometry-predicted damage and erases the Qwen inversion, where CLaRE-ty
states it ``does not establish a formal causal mechanism.''
% tab:clare CONVERTED TO PROSE 2026-07-10 (page budget); all four axes and
% their full content are carried verbatim in the enumeration above.
% Closing positioning paragraph MERGED into the enumeration intro 2026-07-10.
